% !PRINTED 771
\weberpage{771}

In order to prove the irreducibility for an arbitrarily composite $n$,
we apply complete induction. We set
\[
n = p^k n',
\]
and understand by $p$ a prime which divides $n$, but not $n'$, and
assume the \gesperrt{irreducibility of} $f_{n'}(x)$ \gesperrt{to be already proved}. If
then $f_n(x)$ is decomposable into two rational factors
$\varphi(x)$, $\psi(x)$, neither of which is constant:
\weq{7}{f_n(x) = \varphi(x)\, \psi(x),}
then from $\varphi(x)$ we derive the function $\Phi(x)$, whose roots
are the $p^{k}$th powers of the roots of $\varphi(x)$, and which satisfies the condition
\[
\varphi(x) \equiv \Phi(x) \ (\mathrm{mod}\ p)
\]
or also, if we denote by $\chi(x)$ a second function of $x$ with
integral coefficients:
\[
\varphi(x) = \Phi(x) + p\, \chi(x).
\]

But the roots of $\Phi(x)$ are primitive roots of unity of degree
$n'$, and therefore there must be among them at least one,
$\alpha'$, for which $\Phi(\alpha')$ vanishes, so that
\[
\varphi(\alpha') = p\, \chi(\alpha')
\]
holds. On account of the presupposed irreducibility of $f_{n'}(x)$
this equation must, however, hold for \gesperrt{all} roots of unity $\alpha'$, and
for the same reason we have, if $\vartheta(x)$ is a function with
integral coefficients, for all $\alpha'$:
\[
\psi(\alpha') = p\, \vartheta(\alpha').
\]

Accordingly there follows from (7):
\[
f_n(\alpha') = p^2\, \chi(\alpha')\, \vartheta(\alpha'),
\]
and by (4):
\[
1 = p\, \chi(\alpha')\, \vartheta(\alpha')\, g(\alpha').
\]

But the product $\chi(x)\, \vartheta(x)\, g(x)$ is an integral
function of $x$ with integral coefficients, whose remainder on division by $f_{n'}(x)$ likewise has integral coefficients and may be
denoted by $\Theta(x)$. This remainder is of lower degree than $f_{n'}(x)$, and the
function
\[
p\, \Theta(x) - 1
\]
vanishes for \gesperrt{all roots} of $f_{n'}(x)$. Consequently the
equation
\[
p\, \Theta(x) = 1
\]
must be identical, which however is obviously impossible, since $\Theta(x)$ for an
%CONT>>
